% ============================================================
\chapter{Espaces $L^p$}
\label{ch:espaces_lp}
% ============================================================

Les espaces $L^p$ sont les espaces de Banach les plus importants de l'analyse. N\'es de l'int\'egrale de Lebesgue, ils fournissent le cadre naturel pour l'analyse de Fourier, les \'equations aux d\'eriv\'ees partielles, les probabilit\'es et l'apprentissage automatique. Frigyes Riesz, dans les ann\'ees 1910, a \'et\'e le premier \`a \'etudier syst\'ematiquement ces espaces, \'etablissant avec Ernst Fischer le c\'el\`ebre th\'eor\`eme de compl\'etude de $L^2$. Les in\'egalit\'es de H\"older et de Minkowski structurent la th\'eorie~; la dualit\'e $(L^p)^* = L^q$ (pour $1/p + 1/q = 1$) la couronne. Quant \`a l'espace $L^2$, dot\'e de son produit scalaire, il est le seul \`a \^etre un espace de Hilbert, et c'est sur lui que repose toute la th\'eorie spectrale.

\section{D\'efinitions}

\begin{definition}[Espace $L^p$]
Soit $(X, \mathcal{A}, \mu)$ un espace mesur\'e et $1 \leq p < \infty$. L'espace
$\mathcal{L}^p(\mu)$ est l'ensemble des fonctions mesurables $f : X \to \R$
(ou $\C$) telles que~:
\[
\int_X \abs{f}^p\, d\mu < \infty.
\]
La \emph{norme $L^p$} (ou \emph{semi-norme} sur $\mathcal{L}^p$) est~:
\[
\norm{f}_p = \left(\int_X \abs{f}^p\, d\mu\right)^{1/p}.
\]
L'espace $L^p(\mu) = \mathcal{L}^p(\mu) / \sim$ est le quotient par la relation
d'\'equivalence $f \sim g \iff f = g$ $\mu$-p.p.
\end{definition}

\begin{definition}[Espace $L^\infty$]
L'espace $L^\infty(\mu)$ est l'ensemble des classes de fonctions mesurables
essentiellement born\'ees, muni de la norme~:
\[
\norm{f}_\infty = \inf\{M \geq 0 : \abs{f} \leq M \text{ $\mu$-p.p.}\}
= \esssup \abs{f}.
\]
\end{definition}

\begin{remarque}
On travaille avec les classes d'\'equivalence (modulo \'egalit\'e p.p.) pour que
$\norm{\cdot}_p$ soit une vraie norme (et non une semi-norme).
\end{remarque}

\section{In\'egalit\'e de H\"older}

\begin{definition}[Exposants conjugu\'es]
Deux r\'eels $p, q \in [1, +\infty]$ sont \emph{conjugu\'es} (ou \emph{duaux})
si~:
\[
\frac{1}{p} + \frac{1}{q} = 1.
\]
Par convention, $q = \infty$ si $p = 1$ et $q = 1$ si $p = \infty$.
\end{definition}

\begin{lemme}[In\'egalit\'e de Young]
Pour $a, b \geq 0$ et $p, q > 1$ conjugu\'es~:
\[
ab \leq \frac{a^p}{p} + \frac{b^q}{q}.
\]
\end{lemme}

\begin{proof}
La fonction $t \mapsto \log t$ est concave, donc pour $\alpha = 1/p$ et
$\beta = 1/q$~:
\[
\log(\alpha a^p + \beta b^q) \geq \alpha \log(a^p) + \beta \log(b^q)
= \log(a) + \log(b) = \log(ab).
\]
En prenant l'exponentielle~: $ab \leq \frac{a^p}{p} + \frac{b^q}{q}$.
\end{proof}

\begin{theoreme}[In\'egalit\'e de H\"older]
\label{thm:holder}
Soient $p, q \in [1, +\infty]$ conjugu\'es. Si $f \in L^p(\mu)$ et
$g \in L^q(\mu)$, alors $fg \in L^1(\mu)$ et~:
\[
\norm{fg}_1 = \int_X \abs{fg}\, d\mu \leq \norm{f}_p \cdot \norm{g}_q.
\]
\end{theoreme}

\begin{proof}
Le cas $p = 1$, $q = \infty$ est imm\'ediat~: $\abs{fg} \leq \abs{f} \cdot
\norm{g}_\infty$ p.p.

Pour $1 < p < \infty$, on peut supposer $\norm{f}_p = \norm{g}_q = 1$ (sinon
normaliser). Par l'in\'egalit\'e de Young appliqu\'ee \`a $a = \abs{f(x)}$ et
$b = \abs{g(x)}$~:
\[
\abs{f(x) g(x)} \leq \frac{\abs{f(x)}^p}{p} + \frac{\abs{g(x)}^q}{q}.
\]
En int\'egrant~:
\[
\int \abs{fg}\, d\mu \leq \frac{\norm{f}_p^p}{p} + \frac{\norm{g}_q^q}{q}
= \frac{1}{p} + \frac{1}{q} = 1 = \norm{f}_p \cdot \norm{g}_q. \qedhere
\]
\end{proof}

\begin{corollaire}[In\'egalit\'e de Cauchy--Schwarz]
Pour $p = q = 2$~: $\int \abs{fg}\, d\mu \leq \norm{f}_2 \cdot \norm{g}_2$.
\end{corollaire}

\section{In\'egalit\'e de Minkowski}

\begin{theoreme}[In\'egalit\'e de Minkowski]
\label{thm:minkowski}
Soit $1 \leq p \leq \infty$. Pour $f, g \in L^p(\mu)$~:
\[
\norm{f + g}_p \leq \norm{f}_p + \norm{g}_p.
\]
Autrement dit, $\norm{\cdot}_p$ est une norme sur $L^p(\mu)$.
\end{theoreme}

\begin{proof}
Le cas $p = 1$ est l'in\'egalit\'e triangulaire. Le cas $p = \infty$ est \'evident.

Pour $1 < p < \infty$~:
\begin{align*}
\int \abs{f + g}^p\, d\mu &\leq \int \abs{f + g}^{p-1} (\abs{f} + \abs{g})\, d\mu
\\
&= \int \abs{f + g}^{p-1} \abs{f}\, d\mu + \int \abs{f + g}^{p-1} \abs{g}\, d\mu.
\end{align*}
Appliquons H\"older \`a chaque terme avec $q = p/(p-1)$~:
\[
\int \abs{f + g}^{p-1} \abs{f}\, d\mu \leq
\left(\int \abs{f + g}^p\, d\mu\right)^{1/q} \norm{f}_p.
\]
Donc $\norm{f + g}_p^p \leq \norm{f + g}_p^{p/q} (\norm{f}_p + \norm{g}_p)$.
En divisant par $\norm{f + g}_p^{p/q}$ (si non nul) et en notant que
$p - p/q = 1$~: $\norm{f + g}_p \leq \norm{f}_p + \norm{g}_p$.
\end{proof}

\begin{formules}[title=\textbf{In\'egalit\'es fondamentales}]
\[
\text{H\"older : } \norm{fg}_1 \leq \norm{f}_p \norm{g}_q, \qquad
\text{Minkowski : } \norm{f + g}_p \leq \norm{f}_p + \norm{g}_p.
\]
\end{formules}

\section{Compl\'etude : $L^p$ est un espace de Banach}

\begin{theoreme}[Riesz--Fischer]
\label{thm:riesz_fischer}
Pour tout $1 \leq p \leq \infty$, l'espace $(L^p(\mu), \norm{\cdot}_p)$ est un
espace de Banach (espace vectoriel norm\'e complet).
\end{theoreme}

\begin{proof}
Soit $(f_n)$ une suite de Cauchy dans $L^p$. Il suffit de montrer qu'elle poss\`ede
une sous-suite convergente (toute suite de Cauchy ayant une sous-suite convergente
converge).

Choisissons $n_k$ tel que $\norm{f_{n_{k+1}} - f_{n_k}}_p < 2^{-k}$. Posons~:
\[
g_N = \sum_{k=1}^{N} \abs{f_{n_{k+1}} - f_{n_k}}, \qquad
g = \sum_{k=1}^{\infty} \abs{f_{n_{k+1}} - f_{n_k}}.
\]
Par Minkowski, $\norm{g_N}_p \leq \sum_{k=1}^N 2^{-k} \leq 1$. Par convergence
monotone, $\norm{g}_p \leq 1$, donc $g < \infty$ p.p. Cela signifie que la s\'erie
t\'elescopique $f_{n_1} + \sum_{k=1}^{\infty} (f_{n_{k+1}} - f_{n_k})$ converge
absolument p.p. Soit $f$ sa somme (p.p.). On a $f_{n_k} \to f$ p.p. et
$\abs{f_{n_k} - f} \leq 2g \in L^p$, donc par convergence domin\'ee,
$\norm{f_{n_k} - f}_p \to 0$.
\end{proof}

\section{Relations d'inclusion entre espaces $L^p$}

\begin{proposition}
Si $\mu(X) < \infty$ et $1 \leq p \leq q \leq \infty$, alors
$L^q(\mu) \subset L^p(\mu)$ et~:
\[
\norm{f}_p \leq \mu(X)^{1/p - 1/q} \norm{f}_q.
\]
\end{proposition}

\begin{proof}
Appliquer H\"older avec $r = q/p \geq 1$ \`a $\abs{f}^p$ et $1$~:
$\int \abs{f}^p\, d\mu \leq \left(\int \abs{f}^q\, d\mu\right)^{p/q}
\mu(X)^{1 - p/q}$.
\end{proof}

\begin{attention}[title=\textbf{Pas d'inclusion en g\'en\'eral}]
Si $\mu$ n'est pas finie, il n'y a en g\'en\'eral pas d'inclusion entre les $L^p$.
Par exemple, sur $(\R, \lambda)$~: $f(x) = x^{-1/2} \mathbf{1}_{(0,1)} \in L^1
\setminus L^2$ et $g(x) = (1+x)^{-1} \in L^2(\R) \setminus L^1(\R)$.
\end{attention}

\section{Densit\'e et s\'eparabilit\'e}

\begin{theoreme}
Pour $1 \leq p < \infty$~:
\begin{enumerate}
    \item Les fonctions simples int\'egrables sont denses dans $L^p(\mu)$.
    \item Si $X = \R^n$ et $\mu = \lambda^n$, les fonctions continues \`a support
    compact $C_c(\R^n)$ sont denses dans $L^p(\R^n)$.
    \item $L^p(\R^n)$ est s\'eparable pour $1 \leq p < \infty$.
\end{enumerate}
\end{theoreme}

\begin{remarque}
$L^\infty$ n'est en g\'en\'eral \emph{pas} s\'eparable.
\end{remarque}

\section{Dualit\'e des espaces $L^p$}

\begin{theoreme}[Dualit\'e $L^p$--$L^q$]
\label{thm:dualite_lp}
Soit $1 \leq p < \infty$ et $q$ l'exposant conjugu\'e. Pour toute forme lin\'eaire
continue $\ell : L^p(\mu) \to \R$, il existe un unique $g \in L^q(\mu)$ tel que~:
\[
\ell(f) = \int_X fg\, d\mu, \quad \forall f \in L^p(\mu).
\]
De plus, $\norm{\ell} = \norm{g}_q$. Autrement dit, $(L^p(\mu))^* \cong L^q(\mu)$
(isom\'etriquement).
\end{theoreme}

\begin{remarque}
Ce th\'eor\`eme suppose que $\mu$ est $\sigma$-finie. Pour $p = 1$,
$(L^1)^* \cong L^\infty$. Pour $p = \infty$, $(L^\infty)^* \supsetneq L^1$
en g\'en\'eral.
\end{remarque}

\section{L'espace $L^2$ comme espace de Hilbert}

\begin{definition}[Produit scalaire $L^2$]
L'espace $L^2(\mu)$ est muni du produit scalaire~:
\[
\inner{f}{g} = \int_X f \overline{g}\, d\mu.
\]
Il est complet pour la norme associ\'ee $\norm{f}_2 = \sqrt{\inner{f}{f}}$, donc
c'est un espace de Hilbert.
\end{definition}

\begin{theoreme}[Projection orthogonale]
Si $V$ est un sous-espace vectoriel ferm\'e de $L^2(\mu)$ et $f \in L^2(\mu)$, il
existe un unique $\hat{f} \in V$ tel que~:
\[
\norm{f - \hat{f}}_2 = \inf_{g \in V} \norm{f - g}_2.
\]
De plus, $f - \hat{f} \perp V$ (i.e. $\inner{f - \hat{f}}{g} = 0$ pour tout
$g \in V$).
\end{theoreme}

\section{In\'egalit\'es compl\'ementaires}

\begin{proposition}[In\'egalit\'e de Jensen]
\label{prop:jensen}
Soit $(X, \mathcal{A}, \mu)$ un espace de probabilit\'e ($\mu(X) = 1$),
$f \in L^1(\mu)$ et $\varphi : \R \to \R$ convexe. Si $\varphi \circ f$ est
int\'egrable, alors~:
\[
\varphi\!\left(\int_X f\, d\mu\right) \leq \int_X \varphi \circ f\, d\mu.
\]
\end{proposition}

\begin{proposition}[In\'egalit\'e d'interpolation]
Si $f \in L^{p_1} \cap L^{p_2}$ avec $1 \leq p_1 < p_2 \leq \infty$, alors
$f \in L^p$ pour tout $p \in [p_1, p_2]$ et~:
\[
\norm{f}_p \leq \norm{f}_{p_1}^{\theta} \norm{f}_{p_2}^{1-\theta}, \quad
\text{o\`u } \frac{1}{p} = \frac{\theta}{p_1} + \frac{1-\theta}{p_2}.
\]
\end{proposition}

\section{Exercices}

\begin{exercice}
Montrer que pour $1 \leq p < q < \infty$, $\ell^p \subset \ell^q$ (o\`u
$\ell^p = L^p(\N, \mathcal{P}(\N), \#)$ avec $\#$ la mesure de comptage).
C'est l'inverse de l'inclusion pour les espaces \`a mesure finie.
\end{exercice}

\begin{exercice}
Montrer que $L^2([0,1])$ est un espace de Hilbert s\'eparable et que
$(e_n)_{n \in \Z}$ d\'efini par $e_n(x) = e^{2\pi inx}$ est une base hilbertienne.
\end{exercice}

\begin{exercice}
Montrer l'in\'egalit\'e de Minkowski int\'egrale~: si $f \geq 0$ est mesurable sur
$X \times Y$ et $1 \leq p < \infty$, alors
\[
\left(\int_X \left(\int_Y f(x,y)\, d\nu(y)\right)^p d\mu(x)\right)^{1/p}
\leq \int_Y \left(\int_X f(x,y)^p\, d\mu(x)\right)^{1/p} d\nu(y).
\]
\end{exercice}

\begin{exercice}
Soit $(f_n) \subset L^p(\mu)$ avec $\sum_n \norm{f_n}_p < \infty$. Montrer que
$\sum_n f_n$ converge p.p. et dans $L^p$.
\end{exercice}

\begin{exercice}
Montrer que si $f_n \to f$ dans $L^p$ et $f_n \to g$ p.p., alors $f = g$ p.p.
\end{exercice}
